\section{Introduction}

Many datasets record facts that involve more than two things at once. An insurance claim joins a patient, a provider, a drug and a pharmacy. A paper joins its authors, and a meeting joins its participants. A hypergraph stores each such fact as one hyperedge over all of its participants. A property graph has to break the same fact into a new node and one edge per participant, and a relational schema into a join table. Queries then ask hypergraph questions: which facts involve this entity, which facts involve both of these entities, and what lies two hops away.

HyperGraphDB~\cite{hgdb} is the established open-source hypergraph database for the JVM. It stores atoms and links in Berkeley DB Java Edition (JE), a general-purpose log-structured key-value store. HStore is a newer engine built for this data model. Its incidence sets and hyperedge memberships are counted B+trees, and its storage, transaction and recovery design is described in \cref{sec:systems}.

This report compares the two systems at the storage level, on the same machine and data, with the same cache budget. We measure what an embedded application sees: throughput of loads, updates, deletes and reads; per-operation latency, including tail latency while another thread writes; behaviour of giant hyperedges; disk footprint before and after deletes and space reclamation; time to recover after a crash; and the bytes written, heap allocated and GC time behind each result. Every read result is cross-checked between the two systems.

\Cref{sec:method} describes the dataset, workloads and measurement method. \Cref{sec:results} presents the results, and \cref{sec:discussion} explains them in terms of each system's design. \Cref{sec:validity} lists the limits of this study. Full tables, per-run values and the steps to reproduce every number are in the appendices.

\paragraph{Summary of findings.}
\begin{itemize}
\item \textbf{Incidence-based reads} are faster in HStore in every configuration: \M{sumCreadincidence} to \M{sumAreadincidence} for incidence sets, \M{sumCreadcomembership} to \M{sumBreadcomembership} for co-membership counts, and \M{sumCreadincidencecold} to \M{sumAreadincidencecold} directly after a reopen.
\item \textbf{Writes} cost HStore far fewer bytes. Single-update commits write \M{bpoAhstorelatencycommit} against \M{bpoAhgdblatencycommit}, and batched updates run \M{sumBwriteupdate} to \M{sumCwriteupdate} faster. Bulk insertion is closer: between \M{sumBingestedges} and \M{sumCingestedges} for hyperedges.
\item \textbf{Member enumeration} is HStore's weakness once the data outgrows its node cache: \M{sumCreadmembers} at scale 4, and \M{sumCreadtwohop} for two-hop neighbourhoods, which depend on it.
\item \textbf{Tail latency.} HStore's p99 and p99.9 latencies are lower than HyperGraphDB's in every latency workload and configuration, by 1.3 to 69 times. The single slowest commit at scale 4 was HStore's.
\item \textbf{Disk.} HStore's files are \M{sumBdisk} to \M{sumCdisk} smaller after loading. After deletes they first grow while retained history still references the replaced pages.
\item \textbf{Recovery.} Neither system lost an acknowledged commit after SIGKILL. Recovery times overlap.
\end{itemize}

